% called by main.tex
%
\chapter{Aspectos éticos y limitaciones}
\label{ch::etica}

Este capítulo recoge las consideraciones éticas que han guiado el trabajo y una revisión
explícita de sus limitaciones, que delimitan el alcance de cualquier conclusión.

\section{Consideraciones éticas}

El sistema estudia decisiones económicas automatizadas y es especialmente sensible al
sobreajuste cuando el horizonte de datos es reducido. La actuación responsable es no operar
con capital real y no convertir un diagnóstico \emph{post-hoc} en una afirmación de
rentabilidad. La política completa queda congelada para observación prospectiva, sin que los
días usados para formularla cuenten de nuevo como validación.

En cuanto a los datos, se adopta un marco de confidencialidad, integridad y disponibilidad.
Una referencia externa retrasada o errónea puede inducir una señal espuria; por ello el
sistema se entrena solo con sesiones aceptadas, separa datos crudos, características y
etiquetas, mantiene trazabilidad de los artefactos e inhibe la operativa ante datos
obsoletos. El trabajo emplea datos públicos de mercado y no datos personales. El corpus
completo no se distribuye por tamaño; se publican el esquema, un ledger anonimizado de
acciones y el código de evaluación.

\section{Limitaciones}

\textbf{Selección posterior al bloque.} El especialista base sí estaba congelado y su
resultado limpio es $+0{,}349$ ticks por acción al coste de referencia. El filtro de
volatilidad se incorporó después de observar el cambio de régimen; por tanto, su resultado de
$+1{,}069$ ticks sobre el mismo bloque es exploratorio y necesita fechas posteriores.

\textbf{Soporte temporal.} El conjunto fuera de muestra abarca cinco días. Tres de cinco son
positivos para la política base y cinco de cinco dentro del subconjunto diagnóstico. Ninguna
de las dos cifras constituye una validación definitiva.

\textbf{Ausencia del suceso de cola en particiones pequeñas.} La línea de la adenda~D
(Sección~\ref{sec::adenda_d}) reveló una limitación estructural que afecta a cualquier
intento de \emph{aprender} gestión del riesgo con pocos meses de datos: el suceso que la
política debe evitar —el día de pérdida severa— es tan infrecuente (cuatro días de
treinta y cinco en la ventana del corte) que las particiones pequeñas lo pierden con alta
probabilidad. El periodo de entrenamiento de mayo--junio no contiene ni un solo día por
debajo de $-100$\,\$, y un conjunto de validación de una semana sale sin ningún día en
pérdida el 43\,\% de las veces. Cuando eso ocurre, el término de aversión al riesgo de la
función de pérdida vale cero por construcción y el hiperparámetro que la gobierna no puede
seleccionarse: no es un fallo del modelo sino del diseño experimental, y solo lo resuelven
ventanas más largas o la simulación de escenarios de cola.

\textbf{Continuidad de la captura.} El recolector no operó de forma ininterrumpida durante
todo el periodo del proyecto: hay un hueco documentado entre el 14 de junio y la
reanudación real, confirmada el 3 de julio de 2026 a las 08:42 UTC. El hueco no afecta a
ninguna cifra reportada en este trabajo —tanto el corpus principal de entrenamiento
(11-25 de mayo) como el bloque fuera de muestra confirmatorio (6-10 de junio) se
capturaron antes de que se abriera—, pero delimita el origen real de la ventana de
captura continua sobre la que se apoyará la validación prospectiva
(Capítulo~\ref{ch::resultados}).

\textbf{Resolución temporal del instrumento.} La captura opera en una malla de dos segundos,
y esa granularidad limita una de las conclusiones centrales más de lo que su enunciado
sugiere. La latencia de ejecución se midió extremo a extremo contra el mercado real
—429{,}7\,ms en percentil 95 hasta el acuse del CLOB, con un \emph{fill} efectivamente casado
en 426{,}7\,ms— y la decisión operativa registrada fue emplear un segundo como contrato de
modelado, reservando los dos segundos para el escenario de estrés. Sin embargo, un escalón de
un segundo no tiene ninguna observación que lo soporte en una malla de dos, de modo que la
tabla de latencias salta de cero a dos segundos sin ningún punto intermedio. El intervalo
donde vive la latencia realmente medida queda, por tanto, sin observar, y la conclusión
correcta es que no se ha demostrado que la ventaja sobreviva a la latencia de ejecución ni
que deje de hacerlo (Sección~\ref{sec::latencia_grid}). Resolverlo exige capturar a una
frecuencia sensiblemente mayor, con el coste de almacenamiento y de proceso que ello implica;
queda como línea futura. Conviene recordar que la conclusión de fondo se apoya además en un
argumento independiente de la latencia, el de la estructura de comisiones.

\textbf{Cierre de la ventana prospectiva.} La ventana de evaluación prospectiva no la
cerró una decisión metodológica sino un fallo de hardware: el 14 de agosto de 2026 a las
11:05 UTC el disco de almacenamiento del recolector se desconectó del bus USB de la
Raspberry~Pi y la captura se detuvo, reanudándose el 17 de agosto a las 09:23 UTC. El
último día completo capturado es el 13 de agosto, y esa es la fecha de cierre de la
ventana. La integridad de lo capturado se verificó antes de evaluar nada: el sistema de
ficheros se recuperó mediante la reejecución de su diario, el medio no presentó errores
de lectura ni antes ni después del incidente, y la cobertura diaria se comprobó una a una
sobre el registro de sesiones del propio recolector. Resultan 35~días completos e
ininterrumpidos, del 10 de julio al 13 de agosto, por encima del mínimo de 25~días
declarado en el pre-registro. La circunstancia tiene una consecuencia metodológica
favorable que conviene explicitar: el final de la ventana quedó determinado por una causa
externa y ajena al experimento, y quedó determinado \emph{antes} de que se observara
ningún resultado, de modo que la elección del punto de corte no pudo verse influida por
los datos que ese corte iba a revelar. El hueco posterior al 14 de agosto queda
documentado en el monitor del recolector, tal y como exige el propio pre-registro para
cualquier exclusión de días por fallo de captura.

\textbf{Dependencia y solapamiento.} Las 318 unidades del diagnóstico pertenecen a 156
mercados, con hasta ocho señales por mercado y horizontes H60 solapados. El
\emph{bootstrap} por filas subestima esa dependencia. Un remuestreo agrupado por mercado
amplía el IC90\,\% a $[+0{,}193,+2{,}004]$, pero sigue describiendo un subconjunto elegido
\emph{post-hoc}. Además, no se ha fijado una regla de exposición simultánea; por ello la suma
acumulada no es PnL de cartera.

\textbf{Validación \emph{offline}.} La evaluación usa un modelo de ejecución conservador pero
simplificado. No reconstruye la cola orden por orden ni incorpora rellenos parciales. El
registro de sombra actual es retrospectivo; falta una ejecución prospectiva de mayor
fidelidad.

\textbf{Respuesta competitiva del mercado.} Más allá de la fidelidad de la cola, el
simulador maker hereda una limitación estructural de todo \emph{backtest} basado en
reproducción histórica: asume que el resto del mercado —y en particular los creadores de
mercado incumbentes— no reacciona a las órdenes propias. En un régimen de recompensas por
liquidez, esa hipótesis es optimista, porque una orden \emph{maker} nueva compite por la
cola y por el reparto de \emph{rebates} con participantes que ajustarían su comportamiento.
Ningún \emph{replay} puede medir esa reacción; hacerlo exige un simulador reactivo, lo que
se plantea como línea de trabajo futuro (Capítulo~\ref{ch::conclusiones}).

\textbf{Manipulación de \emph{settlement}.} Dai, Jia y Yu~\cite{dai2026settlement}
documentan, tras el lanzamiento del contrato de Bitcoin a cinco minutos de Polymarket,
que el flujo del mercado al contado se dispara en el instante de resolución y provoca
reversiones de precio posteriores, un fenómeno que desaparece en gran medida en los
contratos de quince minutos. Como la vía maker de este trabajo se concentra por tamaño
en los horizontes de una hora y de quince minutos, se auditó si ese patrón aparece en el
corpus propio, comparando la reversión del precio de Bitcoin en los instantes de
\emph{settlement} frente a instantes placebo emparejados por hora. El patrón direccional
aparece con el signo que describe ese trabajo, pero es minúsculo —los contrastes más
fuertes quedan entre $-0{,}2$ y $-0{,}8$ puntos básicos, con $|t|$ entre $2{,}1$ y
$2{,}3$— y ninguno de los seis contrastes sobre el mercado al contado sobrevive la
corrección por comparaciones múltiples de Benjamini-Hochberg al 5\,\%; en los futuros
perpetuos la dirección coincide con estadísticos aún menores. El componente de flujo del
mecanismo no es medible con esta captura —sus volúmenes son agregados móviles de varios
minutos, no cinta instante a instante— y el diseño solo detecta coincidencia temporal,
no atribución causal: un efecto genérico de las marcas redondas de tiempo produciría el
mismo contraste. La auditoría se queda, por tanto, en una comprobación de robustez
favorable a los horizontes empleados, no en una confirmación ni una refutación del
mecanismo.

\textbf{Reproducibilidad parcial.} El repositorio permite recalcular las tablas públicas desde
el ledger anonimizado y probar la lógica compartida. La reproducción integral del
entrenamiento requiere el corpus privado de gran tamaño y se documenta como limitación, no
como capacidad publicada.

\textbf{Patrón en U.} Dentro de la baja volatilidad se observa una posible relación en forma de
U entre volatilidad y resultado. Como se descubrió en los mismos cinco días, se conserva solo
como hipótesis para trabajo futuro.
